% ===== Preambul comun pentru examen, banca de probleme și setul de exerciții (XeLaTeX) =====
% Serii de timp / Time Series Analysis -- Daniel Traian Pele
% Folosire: \documentclass[10.5pt]{article}  \def\examlang{ro} (sau en)  \input{../tsa_exam_preamble}
% Fișierele se compilează din exam/<subfolder>/, deci logo-urile sînt în ../../logos/, iar rezultatele
% software (fișiere text generate de exam/make_figs.py) se includ cu \softout{bank/out/...} sau \softout{practice/out/...}.
\usepackage{fontspec}
\setmainfont{Helvetica Neue}
\setmonofont{Menlo}[Scale=0.92]
\usepackage{amsmath,amssymb}
\usepackage[a4paper,margin=1.9cm,top=2.4cm,headsep=0.5cm,headheight=14pt]{geometry}
\usepackage{graphicx}
\usepackage{float}
\usepackage{booktabs}
\usepackage{array}
\usepackage{enumitem}
\usepackage{xcolor}
\usepackage{fancyhdr}
\usepackage{fancyvrb}
\usepackage{lastpage}
\usepackage[hidelinks]{hyperref}
\definecolor{tsablue}{RGB}{26,58,110}
\definecolor{tsared}{RGB}{205,0,0}
\graphicspath{{../../logos/}{figs/}{../bank/figs/}{../practice/figs/}}

\providecommand{\examlang}{ro}
\newif\ifen
\def\tmpen{en}\ifx\examlang\tmpen\entrue\else\enfalse\fi

% etichete
\ifen
  \renewcommand{\figurename}{Figure}\renewcommand{\tablename}{Table}
  \newcommand{\coursetitle}{Time Series Analysis}
  \newcommand{\subjword}{Problem}
  \newcommand{\rezword}{Solution.}
  \newcommand{\baremword}{Grading:}
  \newcommand{\outword}{Software output}
\else
  \renewcommand{\figurename}{Figura}\renewcommand{\tablename}{Tabelul}
  \newcommand{\coursetitle}{Serii de timp}
  \newcommand{\subjword}{Subiectul}
  \newcommand{\rezword}{Rezolvare.}
  \newcommand{\baremword}{Barem:}
  \newcommand{\outword}{Rezultatul programului}
\fi

% comutator: barem/rezolvare (true) sau subiecte (false)
\newif\ifrez \rezfalse

% antet / subsol
\pagestyle{fancy}
\fancyhf{}
\renewcommand{\headrulewidth}{0.4pt}
\lhead{\small\itshape \coursetitle}
\providecommand{\headyear}{2027}
\rhead{\small\bfseries \headyear}
\cfoot{\small\thepage/\pageref{LastPage}}

\setlength{\parindent}{0pt}
\setlength{\parskip}{0.42ex}
\setlength{\emergencystretch}{3em}

% numerotare automată, continuă, a cerințelor (1, 2, 3, ...)
\newcounter{qq}
\newcommand{\Q}{\refstepcounter{qq}\textbf{\theqq.}~}
% punctaj aliniat la dreapta, pe ultima linie a cerinței
\newcommand{\pts}[1]{\nobreak\hfill\textnormal{\itshape (#1)}\par}

% titlu de subiect: "Subiectul I • Titlu ........ puncte" + linie
\newcommand{\subiect}[3]{%
  \vspace{1.05ex}%
  {\color{tsablue}\bfseries \subjword\ #1 \textbullet\ #2}\hfill{\bfseries #3}\par
  \vspace{0.2ex}{\color{tsablue}\hrule height0.8pt}\vspace{0.7ex}}
% subiect numerotat automat (I, II, ...); argumentul opțional = identificatorul din bancă (vizibil doar în barem)
\newcounter{subj}
\newcommand{\problem}[3][]{%
  \stepcounter{subj}%
  \subiect{\Roman{subj}}{#2\ifrez\if\relax\detokenize{#1}\relax\else\ {\normalfont\footnotesize\color{tsared}[#1]}\fi\fi}{#3}}

% rezultatul unui program (statsmodels, arch, scikit-learn), inclus dintr-un fișier text generat de make_figs.py;
% calea este relativă la exam/ (de exemplu bank/out/ch02_arma.txt)
\newcommand{\softout}[1]{\par\vspace{0.3ex}%
  \VerbatimInput[fontsize=\footnotesize,frame=lines,framesep=1.2mm,rulecolor=\color{tsablue},
    label={\footnotesize\color{tsablue}\outword}]{../#1}\par\vspace{0.2ex}}

% blocul de rezolvare (doar în barem)
\newcommand{\Rez}{\par\vspace{0.2ex}\textit{\rezword}\par\nopagebreak}
\newcommand{\barem}[1]{\par\textit{\baremword\ #1}\par}

% bloc de titlu cu sigle
\newcommand{\titlublock}[1]{%
\begin{center}
  \begin{minipage}[c]{0.16\textwidth}\includegraphics[height=1.15cm]{ase_logo.png}\end{minipage}%
  \begin{minipage}[c]{0.68\textwidth}\centering
  \ifen
    {\bfseries Bucharest University of Economic Studies}\\[0.2ex]
    Faculty of Cybernetics, Statistics and Economic Informatics\\[0.2ex]
    Course: Time Series Analysis
  \else
    {\bfseries Academia de Studii Economice din București}\\[0.2ex]
    Facultatea de Cibernetică, Statistică și Informatică Economică\\[0.2ex]
    Disciplina: Serii de timp
  \fi
  \end{minipage}%
  \begin{minipage}[c]{0.16\textwidth}\raggedleft\includegraphics[height=0.95cm]{ida_logo.png}\end{minipage}\\[1.2ex]
  {\large\bfseries #1}
\end{center}}

% valori utile generale (z_p, chi^2_p = cuantile de ordin p; valori critice Dickey--Fuller pentru eșantioane mari)
\newcommand{\valoriutilegen}{%
\ifen Useful values: $z_{0.975}=1.960$, $z_{0.95}=1.645$, $\chi^2_{0.95}(1)=3.84$, $\chi^2_{0.95}(2)=5.99$,
$\chi^2_{0.95}(3)=7.81$, $\chi^2_{0.95}(4)=9.49$, $\chi^2_{0.95}(5)=11.07$, $\chi^2_{0.95}(6)=12.59$,
$\chi^2_{0.95}(8)=15.51$, $\chi^2_{0.95}(10)=18.31$, $\chi^2_{0.95}(12)=21.03$, $\chi^2_{0.95}(20)=31.41$;
5\% Dickey--Fuller critical values: $-1.95$ (no constant), $-2.86$ (constant), $-3.41$ (constant and trend);
$\ln 2=0.693$ ($z_p$ and $\chi^2_p$ are quantiles of order $p$).
\else Valori utile: $z_{0,975}=1{,}960$; $z_{0,95}=1{,}645$; $\chi^2_{0,95}(1)=3{,}84$; $\chi^2_{0,95}(2)=5{,}99$;
$\chi^2_{0,95}(3)=7{,}81$; $\chi^2_{0,95}(4)=9{,}49$; $\chi^2_{0,95}(5)=11{,}07$; $\chi^2_{0,95}(6)=12{,}59$;
$\chi^2_{0,95}(8)=15{,}51$; $\chi^2_{0,95}(10)=18{,}31$; $\chi^2_{0,95}(12)=21{,}03$; $\chi^2_{0,95}(20)=31{,}41$;
valorile critice Dickey--Fuller la 5\%: $-1{,}95$ (fără constantă); $-2{,}86$ (cu constantă); $-3{,}41$ (constantă
și trend); $\ln 2=0{,}693$ ($z_p$ și $\chi^2_p$ sînt cuantile de ordin $p$).\fi}

% titlul capitolului într-un fișier din bancă (folosit doar ca etichetă de build_variant.py)
\newcommand{\banktitle}[2]{}
